\begin{center}
Problema F

Zé do Grafo

{\tiny Por André da Cunha Ribeiro}

Timelimit: $1$	
\end{center}			

Zé do Grafo é um famoso Mestre Jedi matador de monstros. Diz a lenda que existe uma cidade escondida nos confins do IFGoiano - Campus Rio Verde. A cidade é compostas por várias pontes que ligam vários bairros da cidade, muito parecido com as setes pontes da cidade de Künigsberg. Mas Zé está enfrentando um grande problema na cidade que é o aumento dos monstros.

Mas Zé do Grafo tem contigo um grande poder. Esse poder permitir ele isolar todos os monstros em alguma ilha e depois ele pode destruir as pontes que ligam o bairro. Tudo estaria certo se Zé do Grafo não tivesse o limite de uma bomba. Contudo, ele precisa de sua ajuda para identificar se a cidade possui um bairro ligado por apenas uma ponte.

Ou seja, dadas as descrições dos bairros e das pontes, sua tarefa é determinar se Zé do Grafo consegue separar todos os monstros em um bairro.

\vspace{0.3cm}
\textbf{Entrada}
A entrada contém um caso de teste. A primeira linha do caso de teste contém dois inteiros $N$, $M$, indicando respectivamente o número de bairros $(1 \leq N \leq 50)$ e de pontes $(1 \leq M \leq 200)$.

Cada uma das $M$ linhas seguintes composta de dois inteiros que descreve uma ponte.

\vspace{0.3cm}
\textbf{Saída}
Para o caso de teste da entrada seu programa deve produzir uma linha na saída contendo $'S'$ se Zé do Grafo consegue separar os monstros em um bairro e $'N'$ caso contrário.
\begin{table}[h]
    \centering
    \begin{tabular}{|p{7cm}|p{7cm}|}
        \hline
        \begin{center}
            \vspace{-0.3cm}
            \textbf{Exemplos de Entrada}
            \vspace{-0.3cm}
        \end{center} 
         &  
        \begin{center}
            \vspace{-0.3cm}
            \textbf{Exemplos de Saída}
            \vspace{-0.3cm}
        \end{center} \\ \hline
         \texttt{3 2} & \texttt{S} \\  
         \texttt{1 2} & \texttt{} \\
         \texttt{2 3} & \texttt{} \\
         \hline
         \texttt{6 7} & \texttt{S} \\  
         \texttt{1 2} & \texttt{} \\
         \texttt{1 3} & \texttt{} \\
         \texttt{2 3} & \texttt{} \\
         \texttt{3 4} & \texttt{} \\
         \texttt{4 5} & \texttt{} \\
         \texttt{4 6} & \texttt{} \\
         \texttt{5 6} & \texttt{} \\
         \hline
         \texttt{3 3} & \texttt{N} \\  
         \texttt{1 2} & \texttt{} \\
         \texttt{1 3} & \texttt{} \\
         \texttt{2 3} & \texttt{} \\
         \hline
    \end{tabular}
    \caption{Exemplos de entradas e saídas}
    \label{tab:inputJ}
\end{table}

\newpage

\begin{center}
\noindent
\lstinputlisting{abobora_22/F/ze2022.cpp}
\end{center}
